\section{Implementation Decisions and Trade-offs}

\subsection{Connection reuse vs. a fresh connection per packet}

The HTTP/2 generator opens one \texttt{httpx} client with HTTP/2 enabled and
keeps it open for the whole run, firing several requests at once over that
single connection to actually use multiplexing rather than just claiming to.
The QUIC generator does the same thing for the same reason: a QUIC handshake
costs at least one extra round trip, and if the generator reconnected for
every request that handshake cost would swamp the latency numbers and make
QUIC look artificially slow next to HTTP/2 and MQTT, both of which also keep a
connection open.

The TCP/UDP generator does the opposite on purpose. For every single TCP
packet it opens a new connection, sends, and closes it again. This was not an
oversight. A fresh connect/close cycle produces a real SYN, SYN-ACK and FIN
sequence in the capture, which is exactly what is needed later for the
Failure Visibility analysis and for telling raw TCP apart from the
application-level protocols in a protocol hierarchy view.

\subsection{Mode and pattern as two separate knobs}

The TCP/UDP generator splits "what does a packet look like" from "how is it
paced over time". \texttt{mode} controls size and base interval: normal mode
uses a fixed 512 bytes every 100 ms, stealth mode draws a random size between
64 and 1400 bytes with Poisson-distributed timing. \texttt{pattern}
independently controls the sending cadence on top of that: constant, periodic
burst, or an extra random gap. Keeping these two orthogonal made it possible
to test, for example, stealth-mode sizing combined with a burst cadence, which
would not have been possible if mode and pattern were the same setting.

\subsection{The Poisson pattern on every protocol, not just one}

The assignment only asks for a random, Poisson-distributed pattern somewhere
in the system. This implementation puts it on all four generators using
\texttt{random.expovariate(1/interval)} (or NumPy's exponential distribution
for the TCP/UDP generator), which produces exponentially distributed gaps
between sends with the same average rate as the constant pattern. The reason
to extend it everywhere was to be able to show one single Wireshark capture
where HTTP/2, QUIC and MQTT all use Poisson timing at once, alongside TCP/UDP
stealth mode, rather than only ever demonstrating it on one protocol.

\subsection{Warmup and cooldown double as the ramp pattern}

Rather than building a separate "ramp" mode per protocol, the warmup and
cooldown periods are handled once, centrally, in the controller. A function
called \texttt{\_ramp\_rates()} steps every generator's rate field linearly
from zero up to its configured target over the warmup window, and back down
to zero over the cooldown window, sending the new value through the normal
PATCH endpoint at each step. This single mechanism satisfies both
requirements at once: the warmup/cooldown behavior the assignment asks for,
and the "ramping" sending pattern, since the gradual change shows up as a
slope in a Wireshark I/O graph either way.

\subsection{Adaptive Control}

The controller can run an autonomous control loop, either because a YAML
phase turns it on or because it was switched on manually. Every few seconds it
checks each generator's error rate, and where reported its latency, against
configurable thresholds, and scales the rate up or down by a multiplier
bounded on both ends. This goes past the baseline requirements, but it ties
directly into the resilience patterns covered in the course: the system reacts
to an injected failure on its own, without anyone touching a slider.
